\documentclass{beamer}
\usepackage{beamerthemesplit}
\usepackage{wrapfig}
\usetheme{SPbGU}
\usepackage{pdfpages}
\usepackage{amsmath}
\usepackage{cmap} 
\usepackage[T2A]{fontenc} 
\usepackage[utf8]{inputenc}
\usepackage[english,russian]{babel}
\usepackage{indentfirst}
\usepackage{amsmath}
\usepackage{tikz}
\usepackage{multirow}
\usepackage[noend]{algpseudocode}
\usepackage{algorithm}
\usepackage{algorithmicx}
\usetikzlibrary{shapes,arrows}
\usepackage{fancyvrb}
\usepackage{tikz}
\usepackage{pgfplots}
\usetikzlibrary{calc}
\usetikzlibrary{shapes,arrows}
\usetikzlibrary{arrows,automata}
\usetikzlibrary{positioning}
\usepackage[labelformat=empty]{caption}
\usepackage[labelformat=empty]{subcaption}
\beamertemplatenavigationsymbolsempty


\newcommand{\tikzmark}[1]{\tikz[overlay,remember picture] \node (#1) {};}
\def\Put(#1,#2)#3{\leavevmode\makebox(0,0){\put(#1,#2){#3}}}

\newcommand{\ltz}{$< 1$}
\tikzset{
    state/.style={
           rectangle,
           rounded corners,
           draw=black, very thick,
           minimum height=2em,
           inner sep=2pt,
           text centered,
           },
}

\title[СА + НС]{Комбинирование нейронных сетей и синтаксического анализа для обработки вторичной структуры последовательностей}
\institute[]{
JetBrains Research, Programming Languages and Tools Lab \\
Санкт-Петербургский государственный университет }

\author[Полина Лунина]{Полина Лунина}


\date{14 декабря 2019г.}
\begin{document}
{
\begin{frame}[fragile]
  \begin{tabular}{p{2.0cm} p{7.5cm} p{1cm}}
   \begin{center}
      \includegraphics[height=1.5cm]{pics/jetbrainsResearch.pdf}
    \end{center}
    &
    \begin{center}
      \includegraphics[height=1.5cm]{pics/YC_logo.pdf}
    \end{center}
    &
    \begin{center}
      \includegraphics[height=1.5cm]{pics/SPbGU_Logo.png}
    \end{center}
  \end{tabular}
  \titlepage
\end{frame}
}

\begin{frame} \frametitle{Биоинформатика}
\begin{tabular}{cl}  
    \parbox{0.44\linewidth}{
        \begin{itemize}
            \item Задачи
            \begin{itemize}
                \item Распознавание
                \item Классификация
                \item Предсказание вторичных структур
                \item ...
            \end{itemize}
  \onslide<2,3> {\item Формальное задание вторичной структуры}
  \onslide<3>  {\item Вероятностная оценка}
        \end{itemize}
    }
    & \begin{tabular}{l}
        \vspace{-0.8cm}
        \hspace{-0.8cm}
        \includegraphics[width=5.0cm]{pics/16s.pdf}
    \end{tabular}  \\
\end{tabular}
\end{frame}

\begin{frame}{Наш подход}
\begin{itemize}
    \item Задать основные элементы вторичной структуры (стемы) с помощью грамматики 
    \item Для вероятностной оценки использовать нейронные сети
\end{itemize}

\vspace{1.0cm}
\begin{center}
    \includegraphics[width=8.5cm]{pics/molekula.pdf} 
\end{center}
    
\end{frame}

\begin{frame} \frametitle{Структура решения}
\vspace{0.5cm}
\scriptsize
  \begin{tikzpicture}[->,>=stealth']
  \node[state,
        align=left,
        text width = 3.2cm] (grm)
   {
   \textbf{Грамматика}\\
   Описывает стемы вторичной структуры.
   };

  \node[state,
        below of=grm,
        node distance=1.6cm,
        align = left,
        text width = 3.2cm] (sqs)
   {
   \textbf{Цепочки РНК}\\
   Строки над алфавитом $\{A, C, G, U\}$.
   };

  \node[state,
        below right = -0.4cm and 0.4cm of grm,
        align = left,
        text width = 3.0cm](parser)
  {
  \textbf{Парсер}\\
  Извлекает стемы вторичной структуры из входных данных.
  };

  \node[state,
        right of=parser,
        node distance=4.3cm,
        align = left,
        text width = 4.6cm] (mtrx)
  {
\textbf{Матрицы}\\
{\tiny
\[ \left( \begin{array}{cccc}
0 & 1 & 0 & 1\\
0 & 0 & 1 & 0\\
0 & 0 & 0 & 1\\
0 & 0 & 0 & 0
\end{array} \right)\]}
\vspace{-0.2cm}

  Результат работы парсера для входной строки $\omega$ --- булева матрица $M$, где $ M[i,j] = 1 $ $\iff$  \linebreak  $\text{\ttfamily{s}} \xrightarrow{*}{} \omega[i,j]$.
  };


  \node[state,
        below = 1.2cm of mtrx,
        node distance=3cm,
        align = left,
        text width = 4.6cm](vector)
  {
  \textbf{Входные данные для НС}\\
  Изображения, вектора или другой формат, удобный для конкретной задачи.
  };

  \node[state,
        left = 0.5cm of vector,
        align = left,
        text width = 4.6cm](DNN)
  {
  \textbf{Нейронная сеть}\\

  Решает поставленную задачу распознавания, классификации, генерации и т.д.

  };
  
    \node[state,
        left of = DNN,
        node distance=3.7cm,
        align = left,
        text width = 1.4cm](result)
  {
  \textbf{Результат}
  };

  \path (grm.east) edge (parser.-200)
   (sqs.east) edge (parser.-160)
   (parser) edge (mtrx)
   (mtrx) edge (vector)
   (vector) edge (DNN)
   (DNN) edge (result)
   ;

  \end{tikzpicture}

  \tikzmark{zzz}{
  }
  \pause
  \onslide<2>{\tikz[overlay,remember picture]{\draw[draw=red,thick,double,fill opacity=0.2, rounded corners] ($ (zzz) + (0.11,5.97)$) rectangle ($ (zzz) + (3.46,3.3)$);}}

  \onslide<3>{\tikz[overlay,remember picture]{\draw[draw=red,thick,double,fill opacity=0.2, rounded corners] ($ (zzz) + (3.89,5.25)$) rectangle ($ (zzz) + (7.05,3.865)$);}}

  \onslide<4>{\tikz[overlay,remember picture]{\draw[draw=red,thick,double,fill opacity=0.2, rounded corners] ($ (zzz) + (7.4,6.17)$) rectangle ($ (zzz) + (12.15,2.98)$);}}

  \onslide<5>{\tikz[overlay,remember picture]{\draw[draw=red,thick,double,fill opacity=0.2, rounded corners] ($ (zzz) + (7.4,1.75)$) rectangle ($ (zzz) + (12.15,0.35)$);}}

  \onslide<6>{\tikz[overlay,remember picture]{\draw[draw=red,thick,double,fill opacity=0.2, rounded corners] ($ (zzz) + (2.12,1.75)$) rectangle ($ (zzz) + (6.87,0.33)$);}}

  \onslide<7>{\tikz[overlay,remember picture]{\draw[draw=red,thick,double,fill opacity=0.2, rounded corners] ($ (zzz) + (0.03,1.35)$) rectangle ($ (zzz) + (1.55,0.75)$);}}

\end{frame}

\begin{frame}{Оптимизация}
\textbf{Проблема:} времязатратность синтаксического анализа
\vspace{8pt}

\textbf{Решение:}
    \begin{itemize}
        \item Создать нейронную сеть, обрабатывающую непосредственно цепочку РНК
        \item Обучение в 2 этапа
        \begin{itemize}
            \item Обучить нейронную сеть на результате работы парсера
            \item Расширить ее верхними слоями, принимающими исходную последовательность
        \end{itemize}
    \end{itemize}
\end{frame}

\begin{frame}\frametitle{Эксперименты}
\textbf{Задачи:}
\begin{itemize}
    \item Классификация тРНК эукариотов и прокариотов
    \item Классификация тРНК архей, бактерий, грибов и растений
\end{itemize}
\vspace{6mm}
\textbf{Технологии:}
\begin{itemize}
    \item Платформа YaccConstructor
    \item Библиотека Keras и фреймворк Tensorflow
\end{itemize}
\vspace{6mm}
\textbf{Базы данных:}
\begin{itemize}
    \item tRNADB-CE
    \item Genomic tRNA database
\end{itemize}
\end{frame}

\begin{frame}\frametitle{Эксперименты}
    \begin{itemize}
        \item Два формата представления матриц
        \begin{itemize}
            \item Вектора
            \item Черно-белые изображения
        \end{itemize}
        \item Обучение нейронных сетей на этих данных
        \item Обучение нейронных сетей, принимающих непосредственно последовательности РНК и использующих веса предыдущих моделей
    \end{itemize}
\end{frame}

\begin{frame}{Результаты}
EP --- эукариоты/прокариоты

ABFP --- археи/бактерии/растения/грибы
\vspace{5mm}
\begin{table}[h]
\centering
\caption{Models test results by accuracy metrics}
\begin{tabular}{|l||l|l||l|l|}
\hline
Classifier                                                               & \multicolumn{2}{l||}{EP}               & \multicolumn{2}{l|}{ABFP}           \\ \hline \hline
Approach                                                                 & Vector-based       & Image-based      & Vector-based      & Image-based     \\ \hline
\begin{tabular}[c]{@{}l@{}}Base model\\ accuracy\end{tabular}            & 94.1\%             & 96.2\%           & 86.7\%            & 93.3\%          \\ \hline
\begin{tabular}[c]{@{}l@{}}Extended model \\ accuracy\end{tabular}       & 97.5\%             & 97.8\%           & 96.2\%            & 95.7\%          \\ \hline
\begin{tabular}[c]{@{}l@{}}Total samples\\ (train:valid:test)\end{tabular} & \multicolumn{2}{l||}{20000:5000:10000} & \multicolumn{2}{l|}{8000:1000:3000} \\ \hline
\end{tabular}
\label{acc}
\end{table}
\end{frame}

\begin{frame}{Результаты}
EP --- эукариоты/прокариоты

ABFP --- археи/бактерии/растения/грибы
\vspace{5mm}
\begin{table}[h]
\scriptsize
\centering
\begin{tabular}{|l||l|l|l|l|l|}
\hline
\multirow{2}{*}{Classifier} & \multirow{2}{*}{Class} & \multicolumn{2}{l|}{Vector-based approach} & \multicolumn{2}{l|}{Image-based approach} \\ \cline{3-6} 
                            &                        & precision         & recall        & precision        & recall        \\ \hline \hline
\multirow{2}{*}{EP}         & prokaryotic            & 95.8\%            & 99.4\%        & 96.2\%           & 99.4\%        \\ \cline{2-6} 
                            & eukaryotic             & 99.4\%            & 95.6\%        & 99.4\%           & 99.5\%        \\ \hline \hline
\multirow{4}{*}{ABFP}       & archaeal               & 91.1\%            & 99.2\%        & 91.6\%           & 98.5\%        \\ \cline{2-6} 
                            & bacterial              & 96.6\%            & 95.1\%        & 95.2\%           & 95.5\%        \\ \cline{2-6} 
                            & fungi                  & 98.5\%            & 94.9\%        & 97.5\%           & 94.3\%        \\ \cline{2-6} 
                            & plant                  & 99.4\%            & 95.7\%        & 99.2\%           & 94.7\%        \\ \hline
\end{tabular}
\end{table}
\end{frame}

\begin{frame}{Текущее исследование} 
\begin{itemize}
    \item Парсер находит в цепочке все возможные стемы, однако не все они действительно будут входить в состав вторичной структуры
    \item Хотим сконструировать нейронную сеть, которая отфильтрует лишние контакты между нуклеотидами и предскажет вторичную структуру цепочки
\end{itemize}
\vspace{0.5cm}
\onslide<2>{
\begin{figure}
\centering
\begin{subfigure}{.3\textwidth}
  \centering
  \fbox{\includegraphics[width=.6\linewidth]{pics/structure.png}}
  \caption{Вторичная структура}
\end{subfigure}%
\begin{subfigure}{.3\textwidth}
  \centering
  \fbox{\includegraphics[width=.6\linewidth]{pics/etalon.png}}
  \caption{Contact map}
\end{subfigure}
\begin{subfigure}{.3\textwidth}
  \centering
  \fbox{\includegraphics[width=.6\linewidth]{pics/parsed.png}}
  \caption{Результат парсера}
\end{subfigure}%
\end{figure}
}
\end{frame}

\begin{frame}{Первые эксперименты}
\textbf{Задача:} предсказание вторичных структур цепочек тРНК длины 90

\vspace{6mm}
\textbf{Данные:}
\begin{itemize}
    \item RNAcentral (последовательности тРНК)
    \item CentroidFold (эталонные структуры)
\end{itemize}

\vspace{6mm}

\textbf{Результаты тестирования на 11000 образцов:}
\begin{itemize}
    \item Precision = 84\% (сколько из предсказанных контактов действительно являются контактами в эталоне)
    \item Recall = 89\% (сколько из требуемых контактов было найдено)
\end{itemize}

\end{frame}

\begin{frame}{Планы}
\begin{itemize}
    \item Предсказание вторичных структур для цепочек различных РНК любой длины
    \item Улучшение точности результата путем увеличения количества данных и настройки параметров нейронной сети 
    \item Выбор оптимального источника эталонных данных
    \begin{itemize}
        \item Лучшие результаты на бенчмарках
        \item Возможность предсказания псевдоузлов
    \end{itemize}
\end{itemize}

\end{frame}

\begin{frame}\frametitle{Публикации}
\begin{itemize}
    \item \textbf{BIOINFORMATICS-2019}
      \begin{itemize}
        \item Семён Григорьев, Полина Лунина. The Composition of Dense Neural Networks and Formal Grammars for Secondary Structure Analysis
        \item Публикация: Scopus
      \end{itemize}
      \item \textbf{BIATA-2019} (постерный доклад)
      \begin{itemize}
         \item Семён Григорьев, Полина Лунина. Improved Architecture of Artificial Neural Network for Secondary Structure Analysis
         \item Публикация: BMC Bioinformatics, Scopus
      \end{itemize}
      \item \textbf{CIBB-2019} (доклад)
      \begin{itemize}
        \item Полина Лунина, Семён Григорьев. On Secondary Structure Analysis by Using Formal Grammars and Artificial Neural Networks
        \item Публикация: ожидается
      \end{itemize}
\end{itemize}

\end{frame}


\end{document}
